\section{Abstract Definition Of Boolean Algebra}

Boolean algebra is a mathematical structure that captures the essence of logical reasoning and mirrors the behavior of set operations such as union, intersection, and complement through an algebraic system. Boolean algebra is an algebraic structure:
\[
(B, +, \cdot, ', 0, 1)
\]
where:
\begin{itemize}
\item $B$ is a non-empty set of elements
\item $+$ is a binary operation called \textit{join} (logical OR)
\item $\cdot$ is a binary operation called \textit{meet} (logical AND)
\item $'$ is a unary operation called \textit{complement} (logical NOT)
\item $0$ and $1$ are distinguished elements of $B$, representing the minimum and maximum values, respectively.
\end{itemize}

These components must satisfy the following axioms for all elements $a,b,c \in B$.

\subsection{1. Commutative Law}

This law states that the order of operands does not affect the result:

\[
a + b = b + a
\]

\[
a \cdot b = b \cdot a
\]

\subsection{2. Associative Law}

This law states that grouping of operands does not affect the outcome.

\[
(a + b) + c = a + (b + c)
\]

\[
(a \cdot b) \cdot c = a \cdot (b \cdot c)
\]

\subsection{3. Distributive Law}

This law states that each operation distributes over the other.

\[
a \cdot (b + c) = (a \cdot b) + (a \cdot c)
\]

\[
a + (b \cdot c) = (a + b) \cdot (a + c)
\]

Both distributive laws must hold.

\subsection{4. Identity Law}

This law states that the Boolean algebra contains identity elements for both operations:

\[
a + 0 = a
\]

\[
a \cdot 1 = a
\]

\subsection{5. Complement Law}

This law states that every element has a unique complement such that:

\[
a + a' = 1
\]

\[
a \cdot a' = 0
\]
